\section{Low-frequency-stable mixed formulation and first solver policy}
\label{sec:stable-mixed}

The first implementation freezes a mixed current--potential--charge formulation
that remains finite as $\omega\to0$. It does not eliminate charge by dividing
the continuity equation by $j\omega$.

\subsection{Three-field formulation}

Let $c$ be conductor-current coefficients, $\phi$ scalar-potential
coefficients, and $q$ charge coefficients. Define
\[
  A_c(\omega)=R_c+j\omega L_A(\omega),
\]
where retardation may make $L_A$ frequency dependent. The mixed equations are
\[
  A_c c-D^\ast\phi-\mathcal I^\ast v=0,
  \label{eq:mixed-current}
\]
\[
  Dc+j\omega M_\rho q=0,
  \label{eq:mixed-continuity}
\]
\[
  \phi-\Phi_\rho(\omega)q=0,
  \label{eq:mixed-potential}
\]
\[
  \mathcal I c=i.
  \label{eq:mixed-port}
\]
Together,
\[
 \begin{bmatrix}
  A_c & -D^\ast & 0 & -\mathcal I^\ast\\
  D   & 0        & j\omega M_\rho & 0\\
  0   & I        & -\Phi_\rho & 0\\
  \mathcal I & 0 & 0 & 0
 \end{bmatrix}
 \begin{bmatrix}
  c\\\phi\\q\\v
 \end{bmatrix}
 =
 \begin{bmatrix}
  0\\0\\0\\i
 \end{bmatrix}.
 \label{eq:stable-kkt}
\]

Equation~\eqref{eq:stable-kkt} is the canonical first-implementation
REFERENCE system for a homogeneous conductor scene.

\subsection{DC limit}

At $\omega=0$,
\[
  Dc=0,
\]
while
\[
  A_c(0)c-D^\ast\phi-\mathcal I^\ast v=0.
\]
Thus $\phi$ acts as the finite Lagrange/electric-potential variable enforcing
current continuity. No coefficient proportional to $1/\omega$ is introduced.

The charge variable remains related by
\[
  \phi=\Phi_\rho(0)q
\]
after the gauge is fixed. DC resistance is therefore obtained from the same
mixed system, not from a separate incompatible code path.

\subsection{Why charge elimination is not the primary formulation}

For nonzero frequency, eliminating $q$ gives schematically
\[
  q=-\frac{1}{j\omega}M_\rho^{-1}Dc.
\]
Substitution into the current equation introduces explicit $1/\omega$ scaling.
Although the exact solution may remain finite because $Dc=O(\omega)$, the
algebraic system becomes vulnerable to low-frequency cancellation and poor
conditioning. Such elimination may be used only inside a carefully scaled
preconditioner, not as the canonical physical system.

\subsection{Compatible spaces and discrete inf--sup diagnostic}

Let $M_c$ and $M_\phi$ be positive current and potential/charge norm matrices
after gauge removal. Define the scaled divergence
\[
  \widehat D
  =
  M_\phi^{-1/2}DM_c^{-1/2}.
\]
After excluding the physical solenoidal null space, define
\[
  \beta_h
  =
  \sigma_{\min}^{+}(\widehat D).
\]
A vanishing or rapidly deteriorating $\beta_h$ under basis refinement indicates
an incompatible current/potential pair. The implementation must change or
enrich the basis; it must not hide the defect with an arbitrary diagonal
regularizer.

The charge/potential space is constructed to contain the numerical range of the
current divergence plus required terminal/gauge modes. Range-revealing QR or
SVD may be used during basis construction, with tolerances tied to the
REFERENCE basis error budget.

\subsection{Gauge constraints}

For each electrically isolated component, remove one constant scalar-potential
mode or impose
\[
  g^\ast M_\phi\phi=0
\]
for a chosen constant-mode vector $g$. The corresponding charge neutrality
constraint is imposed consistently. Gauge constraints are assembled separately
from physical terminal-current constraints.

\subsection{Nondimensional block scaling}

Choose diagonal/block scalings
\[
  c=S_c\widetilde c,\quad
  \phi=S_\phi\widetilde\phi,\quad
  q=S_q\widetilde q,\quad
  v=S_v\widetilde v.
\]
$S_c$ is based on conductive-energy/current norms; $S_\phi,S_q$ on electrostatic
energy/capacitance scales; $S_v$ on the port impedance baseline. The scaled KKT
matrix is used for numerical linear algebra, while residual certification is
mapped back to the declared physical Riesz norm.

\subsection{Correctness solver for the first implementation}

The first teacher backend intentionally prioritizes correctness over asymptotic
speed:
\begin{enumerate}
\item assemble dense/high-accuracy interaction blocks for small benchmark
      scenes;
\item apply the gauge projection exactly;
\item solve the scaled KKT system with a pivoted dense factorization;
\item recompute the residual from independently applied physical blocks;
\item perform basis and quadrature enrichment.
\end{enumerate}
This backend establishes trusted truth for small and medium systems and is the
oracle used to validate the iterative backend.

\subsection{Scalable iterative solver}

The scalable REFERENCE backend uses FGMRES or an equivalent robust nonsymmetric
Krylov method on the scaled KKT system. A physics block preconditioner uses
\[
  \widetilde A_c^{-1}
  \approx
  (R_{\rm self}+j\omega L_{\rm self})^{-1}
\]
for conductor blocks and a potential Schur approximation.

If $q$ is formally eliminated only for preconditioning, the potential Schur
operator is
\[
 S_\phi
 \approx
 D A_c^{-1}D^\ast
 +j\omega M_\rho\Phi_\rho^{-1}.
 \label{eq:potential-schur}
\]
The action of $\Phi_\rho^{-1}$ is approximated by intrinsic electrostatic
self-block solves/hierarchical methods; it need not be formed densely.

As $\omega\to0$,
\[
 S_\phi
 \to
 D R_c^{-1}D^\ast,
\]
which supplies the required DC constraint scaling rather than degenerating as
$1/\omega$.

\subsection{Solver acceptance}

An iterative solve is accepted only if:
\begin{enumerate}
\item original mixed-equation Riesz residual passes;
\item port-current constraints pass independently;
\item charge continuity passes independently;
\item gauge residual passes;
\item power closure passes;
\item the result agrees with the dense correctness backend on overlapping
      benchmark sizes.
\end{enumerate}

\subsection{Frequency sweep recycling}

For nearby frequencies or temperatures, previous modal solutions and Krylov
subspaces may be recycled. Recycling changes only the initial subspace; all
acceptance residuals use the current physical operator.

\subsection{Frozen first choice and allowed future alternatives}

Equation~\eqref{eq:stable-kkt} is frozen for the first implementation. Future
loop--tree, quasi-Helmholtz, charge-current, or other stabilized formulations
may replace it only after demonstrating:
\begin{enumerate}
\item the same continuous Scene/port semantics;
\item equal or better low-frequency conditioning;
\item equivalent converged $Z$ and loss observables;
\item no weaker passivity/power certificates.
\end{enumerate}
Thus solver research remains possible without blocking initial code.
